\documentclass{article}
\usepackage[T1]{fontenc}
\usepackage{lmodern}
\usepackage{graphicx}
\usepackage{amsmath,amssymb}
\usepackage[a4paper,margin=2cm]{geometry}
\usepackage[absolute,overlay]{textpos}
\setlength{\TPHorizModule}{1mm}
\setlength{\TPVertModule}{1mm}
\usepackage[indonesian]{babel}
\usepackage{hyperref}
\setlength{\emergencystretch}{2em}
% unit: Halaman 17

\begin{document}

% === Halaman 17 (python/native) ===

17

• Diperoleh pemetaan (cipherteks → plainteks):


P → e


Z → t


W → h


S → a 

• Iterasi 2:

 UZ QSO VUOHXMOPV GPOZPEVSG ZWSZ OPFPESX UDBMETSX AIZ

  t  a         e   e te  a  that  e e a        a    t

 VUEPHZ HMDZSHZO WSFP APPD TSVP QUZW YMXUZUHSX

    e t    ta t  ha e  ee   a e   th     t  a


 EPYEPOPDZSZUFPO MB ZWP FUPZ HMDJ UD TMOHMQ

   e  e e tat  e     the   et

\end{document}
